\documentclass[11pt,a4paper]{article}

\usepackage[utf8]{inputenc}
\usepackage[T1]{fontenc}
\usepackage[brazilian]{babel}
\usepackage{lmodern}
\usepackage[margin=2.5cm]{geometry}
\usepackage{amsmath}
\usepackage{array}
\usepackage{booktabs}
\usepackage{float}
\usepackage{caption}
\usepackage{enumitem}
\usepackage{xcolor}
\usepackage{tikz}
\usepackage{microtype}
\usepackage[hidelinks]{hyperref}

\captionsetup{font=small,labelfont=bf}
\setlist{itemsep=2pt,topsep=4pt}
\setlength{\parskip}{4pt}
\renewcommand{\arraystretch}{1.12}

\newcommand{\evidencia}[1]{\noindent\textbf{Evidência:} \emph{#1}.\par}
\newcommand{\testes}{\noindent\textbf{Testes.} }
\newcommand{\conclusao}{\noindent\textbf{Conclusão.} }

\title{Aprendizado contínuo de detecção de objetos no robô FML Tugger\\[4pt]
\large Relatório técnico dos experimentos}
\author{Araujo\\Dissertação de mestrado, TUM, FML}
\date{6 de novembro de 2026}

\begin{document}
\maketitle

\begin{abstract}
Este relatório descreve um sistema em que um robô logístico aprende classes
novas de objetos durante a operação. O robô captura imagens, o SAM3 encontra a
classe nova, um humano revisa os casos incertos e um único YOLOv8m é retreinado
para aprender a classe sem esquecer as antigas. Um gate estatístico decide se o
modelo vai para o robô, que o troca sem reiniciar. O ciclo completo funcionou
uma vez no robô real, em 02/11/2026: o modelo promovido detecta 7 classes e
perde 2,4\% na média das 5 classes do LOCO. Esse resultado mostra que o sistema
é viável, não que ele é estável: a mesma configuração passou no gate em 2 de 3
seeds numa rodada e em 1 de 3 na repetição. Os principais achados são:
pseudo-rótulos do modelo antigo preservam as classes antigas tão bem quanto ou
melhor que rótulos humanos; a revisão humana não melhorou de forma mensurável a
classe nova; 202 imagens sintéticas elevaram em 12,8\% o AP de um detector
treinado do zero; e a diferença entre ambientes é o maior limite do sistema
(\texttt{iw\_hub}: .968 no laboratório, .277 no conjunto independente).
\end{abstract}

\tableofcontents
\newpage

% ---------------------------------------------------------------------------
\section{Objetivo e convenções}

O objetivo do projeto é permitir que um robô de logística interna passe a
detectar objetos que não existiam no treino original, com pouco trabalho
humano e sem perder o que o detector já sabia. Isso reúne três problemas:
anotar imagens sem um humano desenhar caixas, adicionar classes sem esquecimento
catastrófico e decidir com segurança quando um modelo novo pode ir para o robô.

O ponto de partida é um YOLOv8m~\cite{yolov8} treinado no LOCO~\cite{loco}
com 5 classes: \texttt{small\_load\_carrier} (SLC), \texttt{forklift},
\texttt{pallet}, \texttt{stillage} e \texttt{pallet\_truck}. O sistema adicionou
duas classes: \texttt{iw\_hub}, um robô amarelo do laboratório, e
\texttt{robotic\_arm}, um braço KUKA.

\paragraph{Convenções.} Salvo indicação, os números são AP50 com média de 3
seeds, escritos como no README (.542 significa 0,542). As classes do LOCO são
medidas no LOCO val: 2277 imagens de armazéns nunca vistos no treino. O modelo
de referência (``modelo antigo'') tem SLC .385, forklift .262, pallet .632,
stillage .496 e pallet\_truck .278, com média .411. Cada pergunta da
Seção~\ref{sec:perguntas} indica a força da evidência: \emph{firme} quando o
resultado repete nos seeds ou tem intervalo de confiança claro; \emph{limitada}
quando há amostra pequena, um único seed ou confundidor.

% ---------------------------------------------------------------------------
\section{O sistema}

\subsection{Plataformas}

\begin{itemize}
  \item \textbf{Robô (fml-22):} ROS2 Humble, câmera Azure Kinect e GTX 1060.
    Roda a captura, a sincronização, a revisão e a inferência a 2 Hz.
  \item \textbf{Workstation:} RTX 4090 compartilhada com outros usuários. Roda
    o SAM3, o treino, a avaliação e o gate. Um cron de hora em hora executa
    \path{run_pipeline_cycle.py}.
\end{itemize}

\subsection{Fluxo de dados}

A Figura~\ref{fig:fluxo} mostra o ciclo. Imagens, decisões de revisão e modelos
trafegam por rsync sobre SSH; o ROS fica restrito ao próprio robô.

\begin{figure}[H]
\centering
\begin{tikzpicture}[
  font=\footnotesize, >=stealth,
  box/.style={draw, rounded corners=2pt, align=center, minimum height=9mm, text width=43mm},
  arr/.style={->, thick},
  lab/.style={font=\scriptsize, align=center, fill=white, inner sep=1pt}]
  \node[font=\small\bfseries] at (0,1.1) {Robô};
  \node[font=\small\bfseries] at (9,1.1) {Workstation};
  \node[box] (cap) at (0,0)    {\texttt{capture\_node}\\1 frame a cada 2,5 s, filtro de blur};
  \node[box] (syn) at (0,-1.7) {\texttt{sync\_node}};
  \node[box] (rev) at (0,-5.1) {\texttt{review\_sync\_node}, \texttt{review\_node}\\aprovação no Lichtblick};
  \node[box] (mod) at (0,-8.5) {\texttt{model\_sync\_node}};
  \node[box] (inf) at (0,-10.2) {\texttt{inference\_node}\\2 Hz, GTX 1060};
  \node[box] (ing) at (9,-1.7) {ingestão\\\texttt{dataset/incoming/}};
  \node[box] (sam) at (9,-3.4) {gate do SAM3\\por classe do registry};
  \node[box] (fil) at (9,-5.1) {fila de revisão\\(casos incertos)};
  \node[box] (tre) at (9,-6.8) {lote da receita g e treino\\3 seeds, a partir de 50 aprovações};
  \node[box] (gat) at (9,-8.5) {gate com IC de 95\%\\promoção para produção};
  \draw[arr] (cap) -- (syn);
  \draw[arr] (syn) -- node[lab, above] {rsync a cada 60 s} (ing);
  \draw[arr] (ing) -- (sam);
  \draw[arr] (sam) -- (fil);
  \draw[arr, <->] (fil) -- node[lab, above] {rsync: fila e decisões} (rev);
  \draw[arr] (fil) -- (tre);
  \draw[arr] (sam.east) -- ++(0.5,0) |- node[lab, pos=0.25, right] {auto-\\aprovados} (tre.east);
  \draw[arr] (tre) -- (gat);
  \draw[arr] (gat) -- node[lab, above] {rsync a cada 300 s} (mod);
  \draw[arr] (mod) -- node[lab, right] {serviço de recarga} (inf);
\end{tikzpicture}
\caption{Ciclo de aprendizado entre o robô e a workstation.}
\label{fig:fluxo}
\end{figure}

\subsection{Anotação automática e revisão humana}

Cada classe nova é definida por um prompt de texto em
\path{class_registry.json}. O SAM3~\cite{sam3} roda sobre os frames e usa a
confiança média das caixas: até 0,5 o frame é rejeitado, acima de 0,5 ele
passa, e acima de 0,85 é aprovado automaticamente. Os casos intermediários vão
para a fila de revisão, onde o humano só aprova ou rejeita o frame; as caixas
não são corrigidas. O treino dispara quando há ao menos 50 aprovações novas.

\subsection{Receita de treino}

Cada rodada reinicia do modelo LOCO original, com \texttt{freeze=21}, imgsz 960,
batch 8, 30 épocas fixas, \texttt{last.pt} e 3 seeds. Os pseudo-rótulos usam
sempre confiança 0,5, decisão tomada em 07/10/2026 junto com a de manter um
modelo único. O lote segue a receita g (Tabela~\ref{tab:receita}).

\begin{table}[H]
\centering\small
\caption{Composição do lote de treino (receita g).}
\label{tab:receita}
\begin{tabular}{p{7.2cm}p{7.6cm}}
\toprule
parte & rótulos \\
\midrule
Replay do LOCO, 10 imagens por frame do pipeline & humanos \\
Frames aprovados, até 150 por classe & classes novas: SAM3 e revisão, com as caixas de todas as classes aprovadas no frame; classes do LOCO: modelo antigo \\
Copy-paste das classes novas sobre imagens do LOCO & herdados \\
Frames de ambiente, rejeitados pelo SAM3 para todas as classes & classes do LOCO: modelo antigo \\
Resto do LOCO train & pseudo-rótulos por classe, conforme o registry: SAM3 ou modelo antigo \\
\bottomrule
\end{tabular}
\end{table}

\subsection{Gate de promoção}

O gate (\path{promote_if_better.py}) compara candidato e produção por bootstrap
pareado sobre as mesmas imagens, com 1000 reamostragens. Uma classe antiga só
reprova se o intervalo de 95\% da diferença ficar inteiro abaixo de $-.01$. Uma
classe nova precisa de AP50 de ao menos .10 na primeira promoção. Entre os seeds
aprovados, é promovido o de maior AP50 médio. Antes do treino, um leak check
aborta o lote se ele contiver qualquer imagem de teste.

\subsection{Implantação}

O \texttt{model\_sync\_node} verifica a produção a cada 300 s, baixa o modelo,
troca o arquivo de forma atômica e chama o serviço
\texttt{perception/reload\_model}. O \texttt{inference\_node} recarrega os pesos
sem reiniciar e publica o resultado em \texttt{/perception/detections\_image}.

% ---------------------------------------------------------------------------
\section{Metodologia de avaliação}

\begin{itemize}
  \item \textbf{LOCO:} o split oficial é por armazém. Treino: subsets 2, 3 e 5.
    Val: subsets 1 e 4. O val mede generalização para armazéns nunca vistos.
  \item \textbf{iw\_hub:} teste fixo com 30 frames do laboratório, da mesma
    sessão do treino, e 50 imagens do Roboflow.
  \item \textbf{robotic\_arm:} 84 frames de 05/10 corrigidos à mão, com 109
    braços. Esse dia inteiro fica fora do treino.
  \item \textbf{Ruído:} seeds diferentes variam até cerca de $\pm.05$ nas
    classes fracas. Por isso as comparações usam 3 seeds, batch fixo e IC
    pareado.
\end{itemize}

\paragraph{Bugs de avaliação corrigidos.} A auditoria achou três problemas, e
os números anteriores a ela foram refeitos ou descartados:
\begin{enumerate}
  \item Filtro de confiança 0,5 antes do AP, que cortava a curva
    precisão-recall. O limiar passou a 0,001.
  \item Anotação com id 0, que o COCOeval trata como não casada. Isso derrubava
    \texttt{iw\_hub} no laboratório de cerca de 1,0 para .91--.94.
  \item O split de validação dos fine-tunes continha 170 imagens do LOCO val,
    o que deixava o teste influenciar a escolha do checkpoint. Os experimentos
    limpos usam 30 épocas fixas e \texttt{last.pt}.
\end{enumerate}

% ---------------------------------------------------------------------------
\section{Resultados atuais}

\subsection{Imagens capturadas}

\begin{table}[H]
\centering\small
\caption{Sessões de captura.}
\begin{tabular}{lrp{9.5cm}}
\toprule
sessão & frames & uso \\
\midrule
16/09 & 180 & pipeline \\
17/09 & 340 & pipeline (treino de iw\_hub e do braço; 30 frames viraram teste de iw\_hub) \\
05/10 & 390 & dia reservado para teste: nunca entra no treino \\
02/11, manhã & 172 & pipeline (gravação nova) \\
\textbf{total no pipeline} & \textbf{1082} & \\
02/11, tarde & 391 & fora do pipeline; 88 após tirar duplicatas; sem rótulos \\
\bottomrule
\end{tabular}
\end{table}

\subsection{Do frame capturado ao treino: braço KUKA}

\begin{table}[H]
\centering\small
\caption{Funil dos frames do braço, conferido no banco do pipeline.}
\label{tab:funil}
\begin{tabular}{p{7cm}rr}
\toprule
etapa & com revisão & sem revisão \\
\midrule
frames analisados pelo SAM3 & 1082 & 1082 \\
passaram no filtro (confiança média acima de 0,5) & 221 & 221 \\
depois de tirar o dia de teste (59 frames de 05/10) & 162 & 162 \\
decisão humana (161 revisados, 1 aprovado automaticamente) & 116 aprovados, 46 rejeitados & nenhuma \\
\textbf{no lote, depois de tirar imagens de teste do iw\_hub} & \textbf{107} & \textbf{140} \\
\bottomrule
\end{tabular}
\end{table}

A aprovação humana foi de 72\%. A confiança média do SAM3 no braço ficou entre
0,50 e 0,86, por isso quase nada foi aprovado automaticamente. O lote sem
revisão inclui 33 frames que o humano havia rejeitado.

\subsection{Rodada real do pipeline (02/11)}

A rodada \texttt{pipeline\_finetune\_20261102\_122735} usou 190 frames do robô
(83 de iw\_hub e 107 do braço). Treino dos 3 seeds e gate levaram 54 minutos,
das 12:27 às 13:21. O robô baixou o modelo às 13:43 e passou a detectar 7
classes sem reiniciar. A chamada de recarga estourou o timeout de 15 s do
\texttt{model\_sync\_node}, mas o recarregamento levou cerca de 15,5 s na GTX
1060 e funcionou.

\begin{table}[H]
\centering\small
\caption{Rodada real. Os intervalos são do seed promovido contra a produção.}
\label{tab:rodada}
\begin{tabular}{lrrrrl}
\toprule
classe & antigo & seed 0 (promovido) & seed 1 & seed 2 & IC 95\% do seed 0 \\
\midrule
SLC & .385 & .408 & .403 & .410 & $[+.019,\ +.029]$ \\
forklift & .262 & .225 & .213 & .236 & $[-.078,\ +.010]$ \\
pallet & .632 & .644 & .641 & .644 & $[+.008,\ +.014]$ \\
stillage & .496 & .479 & .470 (reprova) & .481 & $[-.029,\ -.006]$ \\
pallet\_truck & .278 & .251 & .261 & .259 & $[-.044,\ -.0097]$ \\
iw\_hub & --- & .549 & .550 & .535 & \\
braço KUKA & --- & .562 & .524 & .520 & \\
\textbf{média LOCO} & \textbf{.411} & \textbf{.401} & \textbf{.398} & \textbf{.406} & \\
\bottomrule
\end{tabular}
\end{table}

Passaram 2 de 3 seeds. O seed 0 passa em pallet\_truck por margem mínima: o IC
termina em $-.0097$, e o limite é $-.010$. No modelo promovido, \texttt{iw\_hub}
tem .968 nos 30 frames do laboratório e .277 nas 50 imagens do Roboflow. O
número combinado (.549) esconde essa diferença; o que indica generalização é
.277.

\subsection{Pipeline com e sem revisão humana}
\label{sec:revisao}

A mesma rodada foi treinada duas vezes numa cópia do banco. O braço ``com
revisão'' usa os frames aprovados pelo humano; o braço ``sem revisão'' usa todos
os que o SAM3 passou. As caixas são as mesmas nos dois.

\begin{table}[H]
\centering\small
\caption{Com e sem revisão humana. A última coluna indica em quantos seeds o IC
pareado exclui zero.}
\label{tab:revisao}
\begin{tabular}{lrrrl}
\toprule
classe & antigo & com revisão & sem revisão & seeds significativos \\
\midrule
\textbf{braço KUKA} & --- & \textbf{.542} & \textbf{.555} & 0 de 3 \\
iw\_hub (lab e Roboflow) & --- & .542 & .592 & 1 de 3, sem revisão melhor \\
SLC & .385 & .407 & .392 & 3 de 3, sem revisão pior \\
forklift & .262 & .242 & .241 & 0 de 3 \\
pallet & .632 & .644 & .643 & 1 de 3, sem revisão pior \\
stillage & .496 & .476 & .491 & 3 de 3, sem revisão melhor \\
pallet\_truck & .278 & .253 & .236 & 2 de 3, sem revisão pior \\
\textbf{média LOCO} & \textbf{.411} & \textbf{.404} & \textbf{.401} & \\
\bottomrule
\end{tabular}
\end{table}

No braço, os seeds ficaram em .578 / .545 / .503 com revisão e .564 / .547 /
.553 sem revisão. A diferença fica dentro do ruído: a mesma configuração com
revisão, treinada duas vezes, variou até $\pm.02$ por seed (rodada real .562 /
.524 / .520).

\begin{table}[H]
\centering\small
\caption{Resultado do gate contra o modelo LOCO original.}
\label{tab:gate}
\begin{tabular}{lrp{8.4cm}}
\toprule
configuração & aprovados & reprovações (IC 95\% da diferença) \\
\midrule
\textbf{sem revisão} & \textbf{0 de 3} & pallet\_truck: s0 $[-.063,\ -.028]$, s1 $[-.053,\ -.019]$, s2 $[-.063,\ -.026]$ \\
com revisão, rodada real & 2 de 3 & stillage: s1 $[-.037,\ -.015]$ \\
com revisão, repetição & 1 de 3 & stillage: s0 $[-.033,\ -.011]$, s1 $[-.035,\ -.014]$; pallet\_truck: s1 $[-.050,\ -.016]$ \\
\bottomrule
\end{tabular}
\end{table}

\paragraph{Os lotes não diferem só em frames do braço.} A receita mantém
proporções fixas, então os 33 frames a mais mudam todo o lote
(Tabela~\ref{tab:lotes}). O lote sem revisão tem 36\% menos imagens com
pseudo-rótulos do LOCO train, justamente a parte que distingue a receita g, e
mais replay humano. O oráculo (Seção~\ref{q:oraculo}) mostrou que trocar
pseudo-rótulos por rótulos humanos piora pallet\_truck. Logo, a queda de
pallet\_truck sem revisão pode vir dessa mudança de composição, e não dos frames
ruins. O experimento não separa as duas causas.

\begin{table}[H]
\centering\small
\caption{Composição dos dois lotes.}
\label{tab:lotes}
\begin{tabular}{lrr}
\toprule
parte do lote & com revisão & sem revisão \\
\midrule
frames do robô (iw\_hub e braço) & 190 & 223 \\
replay do LOCO (treino / validação) & 1520 / 380 & 1780 / 450 \\
copy-paste & 738 & 867 \\
frames de ambiente & 503 & 503 \\
LOCO train com pseudo-rótulos & 920 & 590 \\
\bottomrule
\end{tabular}
\end{table}

\paragraph{Leitura.} Na classe nova, a revisão não trouxe ganho mensurável.
Nas classes antigas, nenhum modelo sem revisão passou no gate, contra 3 de 6 com
revisão, mas a amostra é pequena e a composição do lote é um confundidor. As 84
imagens de teste do braço foram pré-selecionadas pelo SAM3 e não têm negativos,
o que pode favorecer o treino sem revisão e superestimar o AP nos dois casos.

% ---------------------------------------------------------------------------
\section{Perguntas científicas, testes e conclusões}
\label{sec:perguntas}

\subsection{Qual modelo zero-shot serve como anotador automático?}
\evidencia{limitada}

\testes Grounding DINO tiny~\cite{gdino}, MM-Grounding-DINO
large~\cite{mmgdino}, OWLv2~\cite{owlv2}, LocateAnything-3B
(\texttt{nvidia/LocateAnything}) e SAM3 no LOCO, em amostras de 150 imagens,
com ajuste de prompts e limiares.

\begin{table}[H]
\centering\small
\caption{Detectores zero-shot no LOCO (AP50 geral).}
\begin{tabular}{llr}
\toprule
amostra & modelo e configuração & AP50 \\
\midrule
150 imagens, 4679 caixas & Grounding DINO tiny & .058 \\
& MM-Grounding-DINO large & .000--.008 \\
& OWLv2 base & .125--.143 \\
& SAM3, nomes literais das classes & .195 \\
& SAM3, prompts ajustados & .408 \\
& SAM3, limiar reduzido a 0,2 (muitos falsos positivos) & .475 \\
\midrule
150 imagens, seed 42, 4381 caixas & Grounding DINO tiny & .015 \\
& LocateAnything-3B & .101 \\
& SAM3, prompts ajustados & .378 \\
\midrule
5097 imagens anotadas (train e val) & SAM3, prompts ajustados & .303 \\
\bottomrule
\end{tabular}
\end{table}

\conclusao O SAM3 foi escolhido, com vantagem ampla sobre os outros modelos na
mesma amostra. Nas 5097 imagens, os APs por classe são SLC .369, forklift .261,
pallet .200, stillage .414 e pallet\_truck .272. A comparação não é totalmente
justa: o SAM3 recebeu mais rodadas de ajuste que os outros, e o ensemble de
prompts não superou o melhor prompt único. Com poucas caixas, a amostra pesa:
forklift foi de .631 (14 caixas) a .012 (8 caixas) em amostras de 150 imagens.

\subsection{Qual a qualidade dos rótulos do SAM3, e onde ele falha?}
\evidencia{firme}

\testes Comparação com rótulos humanos na confiança de operação (0,5) e IoU 0,5,
em dois conjuntos do LOCO; \texttt{iw\_hub} em 509 imagens do Roboflow; falsos
alarmes em 500 imagens do LOCO; busca de caixas que juntam pallets.

\begin{table}[H]
\centering\small
\caption{Precisão e recall do SAM3 contra rótulos humanos. Os conjuntos são
diferentes e não devem ser misturados.}
\label{tab:sam3q}
\begin{tabular}{lrrrrrr}
\toprule
& \multicolumn{3}{c}{LOCO train (2538 imagens)} & \multicolumn{3}{c}{LOCO val (2277 imagens)} \\
\cmidrule(lr){2-4}\cmidrule(lr){5-7}
classe & precisão & recall & caixas & precisão & recall & caixas \\
\midrule
SLC & .45 & .36 & 11956 & .59 & .65 & 8911 \\
forklift & .31 & .56 & 438 & .09 & .31 & 124 \\
pallet & .45 & .26 & 61066 & .54 & .32 & 52297 \\
stillage & .58 & .44 & 3149 & .60 & .56 & 2007 \\
pallet\_truck & .57 & .33 & 1531 & .76 & .38 & 1153 \\
\bottomrule
\end{tabular}
\end{table}

\conclusao O SAM3 é excelente em objeto único e distinto: em \texttt{iw\_hub},
precisão .99, recall .73, IoU médio .88 e nenhuma caixa falsa em 500 imagens do
LOCO. No gate, as 509 imagens se dividiriam em 244 aprovadas automaticamente,
128 para revisão e 137 rejeitadas. No braço, contra o teste corrigido, a
precisão é cerca de .98 e o recall cerca de .80. Em cenas densas ele é fraco. Em
pallet, 634 caixas juntam pallets empilhados ou lado a lado, em 459 das 2191
imagens com 4 ou mais pallets (21\%).

\subsection{Um detector treinado sem rótulo humano se aproxima do treinado com rótulos humanos?}
\evidencia{firme}

\testes YOLOv8m treinado a partir de \texttt{yolov8m.pt}, imgsz 960, batch 16,
150 épocas, patience 40 e \texttt{best.pt}, 3 seeds. Mesmas imagens (2538 de
treino, 282 de validação), trocando rótulos humanos por rótulos do SAM3 a 0,5.
Teste no LOCO val com rótulos humanos.

\begin{table}[H]
\centering\small
\caption{Treino do zero com rótulos humanos e do SAM3.}
\begin{tabular}{lrrl}
\toprule
classe & humano & SAM3 & seeds com diferença significativa \\
\midrule
SLC & .349 & \textbf{.427} & 3 de 3, SAM3 melhor \\
forklift & \textbf{.125} & .078 & 2 de 3, SAM3 pior \\
pallet & \textbf{.616} & .319 & 3 de 3, SAM3 pior \\
stillage & .421 & .437 & 0 de 3 \\
pallet\_truck & .173 & \textbf{.213} & 2 de 3, SAM3 melhor \\
\textbf{média} & \textbf{.337} & \textbf{.295} & seeds: .318 / .342 / .351 e .291 / .310 / .283 \\
\bottomrule
\end{tabular}
\end{table}

\conclusao O modelo sem rótulo humano chega a 88\% do modelo com rótulos humanos.
Pallet explica quase toda a perda; sem essa classe, a média é .289 para o SAM3 e
.267 para os humanos. Os modelos humanos ficam abaixo da produção porque a
produção escolheu o \texttt{best.pt} no próprio LOCO val, enquanto aqui a
validação é separada do teste.

\subsection{As técnicas genéricas evitam o esquecimento ao adicionar uma classe?}
\evidencia{limitada}

\testes Congelamento de BatchNorm, replay estratificado, destilação nativa, peso
por classe, copy-paste~\cite{copypaste}, menos camadas congeladas, máscara de
gradiente, treino conjunto e YOLOv8l. Um seed por candidato.

\begin{table}[H]
\centering\small
\caption{Técnicas genéricas (um seed; diferenças abaixo de .05 nas classes fracas
estão dentro do ruído).}
\begin{tabular}{lrrrr}
\toprule
classe & antigo & copy-paste & treino conjunto & YOLOv8l \\
\midrule
SLC & .385 & .376 & .365 & .451 \\
forklift & .262 & .183 ($-30\%$) & .175 ($-33\%$) & .179 ($-32\%$) \\
pallet & .632 & .637 & .631 & .637 \\
stillage & .496 & .459 & .474 & .556 \\
pallet\_truck & .278 & .206 ($-26\%$) & .211 ($-24\%$) & .233 ($-16\%$) \\
iw\_hub (lab / Roboflow) & --- & .997 / .421 & .990 / .338 & 1.000 / .311 \\
\bottomrule
\end{tabular}
\end{table}

\conclusao Nenhuma resolve. A máscara de gradiente piora tudo, porque o
classificador do LOCO se desalinha das features compartilhadas, que continuam
mudando. Até o treino conjunto, com quase todo o LOCO anotado, regride forklift
e pallet\_truck, o que sugere que o problema não é só falta de replay. O YOLOv8l
melhora SLC e stillage, mas é comparado a uma produção em YOLOv8m.

\subsection{Qual receita de modelo único preserva melhor as classes antigas?}
\evidencia{firme}

\testes Oito variantes sob o mesmo protocolo, 3 seeds. Antes disso, rodadas de
self-training na linha de Mayershofer et al.~\cite{selftraining}, com o aluno e
com a produção como professor, e limiares adaptativos na linha de
Song et al.~\cite{apple}.

\begin{table}[H]
\centering\small
\caption{Variantes de fine-tune ao adicionar iw\_hub. A perda é da média LOCO
contra o modelo antigo, recalculada dos arquivos brutos; iw\_hub é o AP50 nas 50 imagens do Roboflow.}
\label{tab:baselines}
\setlength{\tabcolsep}{4pt}
\begin{tabular}{lrrrrrrrr}
\toprule
variante & SLC & fork. & pallet & still. & p.\_truck & média & perda & iw\_hub \\
\midrule
a ingênuo & .150 & .000 & .197 & .156 & .000 & .101 & $-75{,}5\%$ & .363 \\
b replay & .380 & .131 & .639 & .450 & .205 & .361 & $-12{,}0\%$ & .344 \\
c replay e destilação & .379 & .179 & .640 & .456 & .225 & .376 & $-8{,}4\%$ & .344 \\
d copy-paste & .375 & .142 & .638 & .446 & .203 & .361 & $-12{,}1\%$ & .313 \\
e self-training & .372 & .204 & .635 & .451 & .212 & .375 & $-8{,}7\%$ & .250 \\
f banco híbrido & .373 & .230 & .632 & .458 & .240 & .386 & $-5{,}9\%$ & .327 \\
\textbf{g pseudo-rótulos} & .386 & .220 & .636 & .462 & .270 & \textbf{.395} & $\mathbf{-3{,}9\%}$ & .269 \\
h treino conjunto & .373 & .175 & .646 & .465 & .216 & .375 & $-8{,}7\%$ & .374 \\
\bottomrule
\end{tabular}
\end{table}

\conclusao A receita g é a melhor para as classes antigas. Nos 30 frames do
laboratório, \texttt{iw\_hub} fica em .98--1.0 em todas as variantes. Nesse
teste, nenhuma variante de modelo único passou no gate.

As rodadas de self-training mostraram que o professor importa. Com o aluno como
professor, pallet\_truck caiu de .220 para .190 e os rótulos dessa classe
minguaram de 55 para 30 e 19 ao longo das rodadas: o esquecimento se amplifica.
Com a produção como professor, forklift ficou em .256 e pallet\_truck em .216,
e rodadas extras não acrescentaram nada. Os limiares adaptativos subiram
pallet\_truck (.242 para .252), mas derrubaram stillage (.465 para .446) e
forklift (.253 para .214), com média LOCO .386, e foram descartados.

\subsection{O esquecimento vem de rótulos ruins?}
\label{q:oraculo}
\evidencia{firme}

\testes Oráculo: mesmas 1760 imagens do LOCO train e mesma receita, trocando
pseudo-rótulos por rótulos humanos em stillage e pallet\_truck, e depois em todas
as classes. 3 seeds.

\begin{table}[H]
\centering\small
\caption{Anotador oráculo.}
\begin{tabular}{lrrrrrr}
\toprule
braço & SLC & forklift & pallet & stillage & pallet\_truck & média \\
\midrule
pseudo-rótulos & .391 & .233 & .637 & .469 & \textbf{.261} & \textbf{.398} \\
humanos em stillage e pallet\_truck & .393 & .254 & .637 & .465 & .222 & .394 \\
humanos em todas as classes & .368 & .199 & .645 & .474 & .221 & .382 \\
\bottomrule
\end{tabular}
\end{table}

\conclusao Não. Rótulos humanos não reduzem o esquecimento, e pallet\_truck
piora mesmo com mais caixas (1002 contra 589 no lote). A interpretação proposta
é que pseudo-rótulos do modelo antigo agem como destilação~\cite{hinton}: o
modelo novo imita o antigo, inclusive nos erros. É o achado mais original do
trabalho, mas o mecanismo foi inferido, não testado diretamente.

\subsection{Isolar a classe nova num segundo modelo evita o esquecimento?}
\evidencia{firme}

\testes Modelo LOCO intacto mais um especialista YOLOv8s, no estilo de isolamento
de parâmetros~\cite{progressive}, com gate de falsos positivos de até 1\% das
imagens do LOCO. 3 seeds.

\conclusao Evita por construção: as classes do LOCO ficam idênticas. O custo é de
57 ms por frame no robô, contra 38 ms do modelo sozinho. Passa em 2 de 3 seeds,
com taxas de falso positivo de 0,18\%, 2,06\% e 0,35\%; o seed 1 reprova. A
abordagem foi descartada em 07/10 pela decisão de manter um modelo único.

\subsection{Até quantas classes o modelo único aguenta?}
\evidencia{limitada}

\testes Adição incremental de até 70 classes do COCO~\cite{coco} (sem
\texttt{person}), com 300 imagens de treino e 200 de avaliação por classe.

\begin{table}[H]
\centering\small
\caption{AP50 médio por grupo fixo de classes no teste de saturação.}
\begin{tabular}{lrrrrr}
\toprule
grupo & logo após aprender & 20 classes & 40 classes & 70 classes & perda \\
\midrule
1--10 & .112 & .098 & .082 & .079 & $-29\%$ \\
11--20 & .123 & --- & .116 & .103 & $-16\%$ \\
21--40 & .078 & --- & --- & .072 & $-7\%$ \\
\bottomrule
\end{tabular}
\end{table}

\conclusao A queda do primeiro grupo quase para entre 40 e 70 classes. O LOCO cai
de .411 para .364 com 40 classes ($-11\%$). O AP absoluto é baixo (11 das 70
classes abaixo de .01), então o teste mede tendência, não um modelo bom em cada
classe.

\subsection{Como decidir com confiabilidade se um modelo pode ser implantado?}
\evidencia{firme}

\testes Auditoria da avaliação (Seção~3); gate com bootstrap pareado e IC de 95\%;
leak check.

\conclusao Uma tolerância fixa de $-.01$ é menor que o ruído entre seeds; o IC
resolve isso. Dois vazamentos reais apareceram:
\begin{itemize}
  \item Um banco de imagens sorteado nos 5 armazéns, incluindo os 2 de validação,
    deu média .491, um falso +19\%. Refeito só com os armazéns de treino, o mesmo
    método deu .393. Esse caso motivou o leak check.
  \item Na rodada real, 9 frames aprovados para o braço eram imagens de teste de
    iw\_hub. O leak check barrou a primeira tentativa, e imagens de teste de
    qualquer classe passaram a ser excluídas.
\end{itemize}

\subsection{O sistema aprende uma segunda classe sem esquecer a primeira?}
\evidencia{limitada}

\testes Braço KUKA, 3 seeds, com o mesmo teste de 84 imagens. Primeiro offline,
a partir do seed 2 da receita g e sem revisão (120 frames com braço, 199 caixas).
Depois, na rodada real, a partir do modelo LOCO original.

\begin{table}[H]
\centering\small
\caption{Segunda classe: experimento offline e rodada real.}
\begin{tabular}{lrrr}
\toprule
classe & base (g, seed 2) & offline, com braço & rodada real (média de 3 seeds) \\
\midrule
SLC & .384 & .399 & .407 \\
forklift & .229 & .226 & .225 \\
pallet & .636 & .635 & .643 \\
stillage & .463 & .462 & .477 \\
pallet\_truck & .272 & .237 & .257 \\
iw\_hub & .531 & .584 & .545 \\
braço KUKA & --- & .305 & .535 \\
gate & --- & 0 de 3 & 2 de 3 \\
\bottomrule
\end{tabular}
\end{table}

\conclusao O braço foi de .305 para .52--.56. As duas configurações mudam em
várias coisas: modelo de partida, revisão, gravação nova e fusão de rótulos por
frame. Por isso a causa do ganho não foi isolada, e o gate offline usou outra
referência (a própria base). O teste vem do mesmo ambiente do treino e foi
pré-selecionado pelo SAM3.

\subsection{O ciclo completo funciona de ponta a ponta no robô real?}
\evidencia{limitada}

\testes Rodada real em 02/11: gravação por teleoperação (172 frames, das 10:45
às 10:52), SAM3, revisão, treino, gate, promoção, download pelo robô e inferência
ao vivo. A demonstração foi gravada em vídeo.

\conclusao Sim, uma vez. Foi a primeira promoção desde o modelo inicial, com
margem mínima em pallet\_truck e 2 de 3 seeds; a repetição da mesma configuração
passou em 1 de 3. Isso mostra viabilidade, não estabilidade ao longo de várias
rodadas.

\subsection{Quanto a revisão humana agrega para uma classe nova?}
\evidencia{limitada}

\testes Mesma rodada com os 107 frames aprovados e com os 140 que o SAM3 passou,
3 seeds, IC pareado (Seção~\ref{sec:revisao}).

\conclusao Nenhum ganho mensurável no braço (.542 contra .555, 0 de 3 seeds). Sem
revisão, 0 de 3 modelos passam no gate, todos por pallet\_truck, contra 3 de 6 com
revisão. A revisão pode proteger as classes antigas, mas a amostra é pequena e
os lotes diferem também na composição, inclusive em 920 contra 590 imagens com
pseudo-rótulos. O resultado vale para revisão que só filtra frames e para este
teste, pré-selecionado pelo SAM3.

\subsection{Imagens sintéticas melhoram a detecção em imagens reais?}
\evidencia{firme}

\testes Rótulos humanos do LOCO mais 202 renders do Isaac Sim~\cite{isaac} (7,4\%
das 2740 imagens de treino), 3 seeds, validação e teste só com imagens reais. Os
renders têm 97 caixas de SLC, 123 de forklift, 70 de pallet, 94 de stillage e 104
de pallet\_truck. Uma primeira rodada perdeu os rótulos de pallet por um erro no
mapeamento de classes e foi refeita; ela havia dado média .362.

\begin{table}[H]
\centering\small
\caption{Treino do zero com e sem imagens sintéticas.}
\begin{tabular}{lrrrl}
\toprule
classe & humano & com sintético & ganho & seeds com diferença significativa \\
\midrule
SLC & .349 & .380 & $+9\%$ & 2 de 3 melhor, 1 sem diferença \\
forklift & .125 & \textbf{.221} & $+77\%$ & 2 de 3 melhor, 1 sem diferença \\
pallet & .616 & .622 & $+1\%$ & 2 de 3 melhor, 1 de 3 pior \\
stillage & .421 & .463 & $+10\%$ & 3 de 3 melhor \\
pallet\_truck & .173 & .216 & $+25\%$ & 3 de 3 melhor \\
\textbf{média} & \textbf{.337} & \textbf{.380} & $+12{,}8\%$ & seeds: .362 / .384 / .395 \\
\bottomrule
\end{tabular}
\end{table}

\conclusao Sim, nenhuma classe piora de forma consistente. O maior ganho vem da
classe mais rara: forklift tem 438 caixas reais, e o sintético adiciona 123. Vale
para treino do zero e uma única proporção. Forklift varia muito entre seeds (.142
a .309). Dentro do pipeline contínuo, isso ainda não foi testado.

\subsection{O aprendido no robô generaliza para fora do ambiente?}
\evidencia{firme}

\testes \texttt{iw\_hub} em 30 frames do laboratório contra 50 imagens do Roboflow;
banco de imagens de armazéns de teste contra armazéns de treino; val contra teste
no treino do zero.

\conclusao Pouco. \texttt{iw\_hub} cai de .968 para .277. O mesmo método dá .491
com imagens dos armazéns de teste e .393 com outros armazéns. No treino do zero,
o mAP50 cai de cerca de .77 no val para cerca de .30--.38 no teste. A diferença
entre ambientes é o principal limite do sistema.

\subsection{Um modelo de linguagem com visão substitui o anotador?}
\evidencia{limitada}

\testes Claude Opus 5.5 anotando à mão 26 imagens do LOCO val, com grade de
coordenadas, contra rótulos humanos a IoU 0,5. O teste foi interrompido antes das
100 imagens planejadas.

\begin{table}[H]
\centering\small
\caption{Modelo de linguagem com visão como anotador.}
\begin{tabular}{lrrrrr}
\toprule
classe & caixas humanas & caixas do modelo & acertos & precisão & recall \\
\midrule
pallet & 889 & 198 & 76 & .38 & .09 \\
pallet\_truck & 5 & 6 & 4 & .67 & .80 \\
\bottomrule
\end{tabular}
\end{table}

\conclusao Não para cenas densas. O LOCO tem cerca de 34 pallets por imagem, 71\%
deles em prateleiras altas, e quase nenhum destes foi anotado. Das 198 caixas, 65
ficaram com IoU entre 0,3 e 0,5: objeto certo, caixa imprecisa. Pallet\_truck vai
bem, mas com só 5 caixas. As 6 imagens usadas para aprender as convenções estão
incluídas nesses números. O uso como revisor de caixas não foi testado.

\subsection{A receita funciona em um benchmark externo?}
\evidencia{limitada}

\testes CUB~\cite{cub}, protocolo de Incremental Zero-Shot Learning de Wei et
al.~\cite{izsl}: split padrão de ZSL, 5 etapas de 30 classes, SAM3 para as 50
classes não vistas.

\begin{table}[H]
\centering\small
\caption{CUB: classes não vistas (u), vistas (s) e média harmônica (H).}
\begin{tabular}{llrrr}
\toprule
método & guarda imagens antigas? & u & s & H \\
\midrule
Wei et al. & não & 38,1 & 45,6 & 41,5 \\
nosso, sem replay & não & 33,4 & 13,7 & 19,5 \\
nosso, 5 imagens por classe & 5 por classe & 44,7 & 26,0 & 32,9 \\
nosso, replay completo & todas & 44,1 & 46,3 & 45,2 \\
treino conjunto (teto do paper) & todas & 43,7 & 57,7 & 49,7 \\
\bottomrule
\end{tabular}
\end{table}

\conclusao Com replay completo, H fica acima do paper; sem replay, o método
colapsa. Ele depende de guardar dados antigos, o que o paper não faz, e o
pré-treino do SAM3 pode incluir essas espécies. Serve para o cenário em que o
dataset antigo continua disponível.

% ---------------------------------------------------------------------------
\section{Limitações e pendências}

\begin{itemize}
  \item Só 3 seeds, um único ambiente de robô e uma única rodada real.
  \item A aprovação no gate está perto do limite e varia entre repetições.
  \item Teste do braço: 84 imagens do mesmo ambiente, pré-selecionadas pelo SAM3
    e sem negativos.
  \item Teste do \texttt{iw\_hub}: só 50 imagens do Roboflow. É um conjunto
    público de terceiros, de 2024, e a origem precisa ser confirmada
    (Apêndice~\ref{ap:notas}). A gravação de 02/11 à tarde segue sem rótulos.
  \item Muitos frames quase duplicados: na gravação da tarde, 391 viraram 88.
  \item A ablação da revisão não isola a qualidade dos frames: a composição do
    lote também muda.
  \item Mecanismo de destilação inferido, não testado. A proteção de pallet\_truck
    continua em aberto.
  \item Dados sintéticos ainda não foram testados dentro do pipeline.
  \item Regras da TUM sobre declarar uso de IA: conferir antes da entrega.
\end{itemize}

% ---------------------------------------------------------------------------
\appendix

\section{Notas de conferência dos dados}
\label{ap:notas}

Ao preparar este relatório, os números foram conferidos nos arquivos brutos e no
banco do pipeline. Estes pontos diferem do README ou do registro de experimentos,
e aqui vale a versão abaixo.

\begin{enumerate}
  \item \textbf{SAM3 ``no val completo'' (.30).} O valor .303 vem de 5097 imagens,
    o que corresponde a treino (2820) e val (2277) somados, não só ao val.
  \item \textbf{Diferença em \path{sam3_label_quality.json}.} Esse arquivo mede o
    LOCO val; a tabela do README mede o LOCO train. São conjuntos diferentes, e a
    Tabela~\ref{tab:sam3q} traz os dois. A causa exata da diferença não foi
    investigada.
  \item \textbf{Frames retidos de 05/10.} O banco tem 59 frames do braço marcados
    como \texttt{held\_out\_test}, o que fecha a conta 221 $-$ 59 = 162. O registro
    de experimentos dizia 74.
  \item \textbf{Perda da receita g.} Recalculada, é $-3{,}9\%$, não $-3{,}8\%$.
  \item \textbf{Confundidor da revisão.} O README cita a diferença no número de
    frames, mas não a composição do lote (Tabela~\ref{tab:lotes}).
  \item \textbf{Forklift do SAM3 (.012).} Vem de uma amostra com só 8 caixas; nas
    5097 imagens o valor é .261.
  \item \textbf{Origem do Roboflow.} O README trata o conjunto como independente,
    mas o docstring de \path{sam3_label_quality.py} diz que ele vem do laboratório
    do robô. A página do dataset não confirma nenhum dos dois. Se vier do mesmo
    laboratório, .277 mede outra sessão, não outro ambiente.
  \item \textbf{README do pipeline.} \path{dataset/pipeline/README.md} ainda
    descreve o banco de 5000 imagens e diz que as rodadas não promovem; os dois
    pontos estão desatualizados.
\end{enumerate}

\section{Onde estão os dados}

\begin{table}[H]
\centering\small
\caption{Resultados e scripts de cada experimento, em \texttt{dataset/pipeline/}.}
\begin{tabular}{p{3.6cm}p{4.6cm}p{6.6cm}}
\toprule
experimento & script & resultados \\
\midrule
Variantes de receita & \path{baselines.py} & \path{self_training/results_baselines.json} \\
Oráculo & \path{oracle_annotator.py} & \path{self_training/results_oracle_annotator.json} \\
Qualidade do SAM3 & \path{sam3_label_quality.py} & \path{sam3_label_quality.json} \\
SAM3 contra humano & \path{loco_sam3_vs_human.py} & \path{self_training/results_loco_sam3_vs_human.json} \\
Sintético & \path{loco_synthetic_mix.py} & \path{self_training/results_loco_synthetic_mix.json} \\
Braço offline & \path{kuka_arm_experiment.py} & \path{self_training/results_kuka_arm.json} \\
Rodada real & \path{run_pipeline_cycle.py} & \path{self_training/results_pipeline_round_20261102.json} \\
Com e sem revisão & \path{kuka_review_ablation.py} & \path{self_training/results_kuka_review_ablation.json} e \path{_gate.json} \\
Detectores zero-shot & \path{eval_loco.py} (em \texttt{dataset/}) & \path{dataset/results/loco_eval_results_zeroshot.csv} \\
\bottomrule
\end{tabular}
\end{table}

\begin{thebibliography}{99}
\bibitem{loco} C.~Mayershofer, D.-M.~Holm, B.~Molter e J.~Fottner.
  LOCO: Logistics Objects in Context. \emph{IEEE ICMLA}, 2020.
\bibitem{yolov8} G.~Jocher, A.~Chaurasia e J.~Qiu. Ultralytics YOLOv8. Software, 2023.
\bibitem{sam3} N.~Carion et al. SAM 3: Segment Anything with Concepts. 2025.
\bibitem{gdino} S.~Liu et al. Grounding DINO: Marrying DINO with Grounded
  Pre-Training for Open-Set Object Detection. \emph{ECCV}, 2024.
\bibitem{mmgdino} X.~Zhao et al. An Open and Comprehensive Pipeline for Unified
  Object Grounding and Detection. 2024.
\bibitem{owlv2} M.~Minderer, A.~Gritsenko e N.~Houlsby. Scaling Open-Vocabulary
  Object Detection. \emph{NeurIPS}, 2023.
\bibitem{selftraining} C.~Mayershofer, T.~Fischer e J.~Fottner. Improving visual
  object detection using synthetic self-training. \emph{Logistics Journal}, 2021.
\bibitem{apple} Song et al. Overcoming Catastrophic Forgetting for Multi-Label
  Class-Incremental Learning. \emph{WACV}, 2024.
\bibitem{izsl} K.~Wei, C.~Deng, X.~Yang e D.~Tao. Incremental Zero-Shot Learning.
  \emph{IEEE Transactions on Cybernetics}, 2022.
\bibitem{hinton} G.~Hinton, O.~Vinyals e J.~Dean. Distilling the Knowledge in a
  Neural Network. 2015.
\bibitem{progressive} A.~Rusu et al. Progressive Neural Networks. 2016.
\bibitem{copypaste} G.~Ghiasi et al. Simple Copy-Paste is a Strong Data
  Augmentation Method for Instance Segmentation. \emph{CVPR}, 2021.
\bibitem{coco} T.-Y.~Lin et al. Microsoft COCO: Common Objects in Context.
  \emph{ECCV}, 2014.
\bibitem{cub} C.~Wah et al. The Caltech-UCSD Birds-200-2011 Dataset. Relatório
  técnico, Caltech, 2011.
\bibitem{isaac} NVIDIA. Isaac Sim. Software.
\end{thebibliography}

\end{document}
